\documentclass[11pt]{article}
\usepackage[margin=1in]{geometry}
\usepackage{amsmath,amssymb,amsthm,mathtools,microtype}
\usepackage[T1]{fontenc}
\usepackage{lmodern}
\usepackage[colorlinks=true,linkcolor=blue,citecolor=blue,urlcolor=blue]{hyperref}
\newtheorem{theorem}{Theorem}
\newtheorem{lemma}[theorem]{Lemma}
\newtheorem{proposition}[theorem]{Proposition}
\theoremstyle{remark}\newtheorem{remark}[theorem]{Remark}
\newcommand{\R}{\mathbb R}\newcommand{\C}{\mathbb C}
\newcommand{\ip}[2]{\langle #1,#2\rangle}
\newcommand{\norm}[1]{\lVert #1\rVert}
\newcommand{\epi}{\operatorname{epi}}
\newcommand{\dist}{\operatorname{dist}}
\newcommand{\Rea}{\operatorname{Re}}
\newcommand{\Ima}{\operatorname{Im}}
\title{Nonunique nearest stable complex quadratics}
\author{Alec Kriebel\\\small Independent researcher\\
\small ORCID: \href{https://orcid.org/0009-0001-9320-500X}{0009-0001-9320-500X}}
\date{11 October 2026\quad\textnormal{(research note; unrefereed preprint)}}
\begin{document}
\maketitle
\begin{abstract}
The closed set of complex monic quadratics with roots in the closed left half-plane
is not prox-regular at $z^2$. For every fixed positive-definite real inner product
on coefficient space, unstable quadratics with multiple attained global nearest
stable quadratics occur arbitrarily close to $z^2$. The conclusion also holds in
the full exact-degree quadratic coefficient space. A direct proof exhibits a
limiting normal whose cubic support dominates a quartic squared displacement.
The same obstruction disproves set prox-regularity of the whole abscissa epigraph.
In contrast, the monic quadratic abscissa is function prox-regular at the origin
relative to each finite limiting subgradient. Thus the historical question about
local unique stable-polynomial projection has a negative answer for complex
coefficients, while this example does not refute the separate function condition.
\end{abstract}

\section{Question, conventions, and result}
Lewis \cite[\S8, pp.~166--168]{Lewis2003} asks whether the stable monic set is
prox-regular, after specifying complex coefficients and their real Euclidean
coefficient metric. The later Oberwolfach question \cite[p.~106, Question~1]{OWR}
asks whether every degree-$k$ polynomial near $z^k$ has a unique nearest stable
polynomial. That short question does not specify its coefficient field,
normalization, or metric. We address its complex coefficient convention, which
is explicit in the earlier set question and in the cited variational theory
\cite{BO,BLO}. No conclusion for the real-only general question is asserted.

For $p_{a,b}(z)=z^2+az+b$, identify $(a,b)\in\C^2$ with $\R^4$, in the order
$(\Rea a,\Ima a,\Rea b,\Ima b)$. Let $\alpha(p)$ be the maximum real part of
a root, and set
\[
 S=\{(a,b):\alpha(p_{a,b})\leq0\},\qquad c_0=(0,0).
\]
Imaginary-axis roots, including multiple roots, are allowed. For a fixed real
symmetric positive-definite matrix $G$, write
$\ip{u}{v}_G=u^TGv$ and $\norm{u}_G^2=u^TGu$.
In the full coefficient space put
\[
 S_{\rm full}=\{(\ell,a,b)\in\C^3:\ell\ne0,
                 \alpha(\ell z^2+az+b)\leq0\},
 \qquad c_{\rm full}=(1,0,0).
\]
A nearest point always refers to a global minimum over the stated stable set.

\begin{theorem}\label{thm:main}
For every fixed positive-definite real coefficient inner product, $S$ is not
prox-regular at $c_0$, and every neighborhood of $c_0$ contains an unstable
quadratic with at least two distinct attained global nearest points in $S$.
The corresponding statements hold for $S_{\rm full}$ at $c_{\rm full}$.
The abscissa epigraph is not prox-regular as a set at height zero over either
base point, for every fixed positive-definite real inner product on the
respective coefficient/height space.
\end{theorem}

The first-order polynomial root and normal calculus and the connection between
prox-regularity and projection are classical \cite{BO,Lewis2003,PRT}.
The underlying stability-boundary cusp is a positive Whitney-umbrella geometry;
the related real-quartic boundary was studied by Bottema \cite{Bottema}
and is reviewed as a Whitney umbrella by Kirillov--Overton \cite{KO}.
We give the coefficient normal, metric obstruction, and global
attainment arguments directly. This note makes no assertion of firstness or
of the question's present-day openness. The accompanying priority memorandum
records the primary sources checked and the remaining inaccessible bodies.
Nearest-stable coefficient approximation was already studied by Moses--Liu
\cite{ML} for real disk-stable polynomials, including recognition of possible
global nonuniqueness, and by Orban de Xivry \cite[Chapter~4]{Orban} for complex
half-plane coefficients. Our assertion concerns attained ties in every
neighborhood of the specified quadratic base, with all fixed coefficient
inner products; generic nonconvexity or competing local algorithmic minima
alone do not establish that assertion.

\section{A limiting normal that is not proximal}
For a locally closed set $C$ in a finite-dimensional inner product space,
a vector $v$ is a proximal normal at $x\in C$ if, for some $\tau>0$, $x$ is a
nearest point to $x+\tau v$; equivalently, a local quadratic upper support
can be used and the normal step made sufficiently small. Limits of proximal
normals at convergent points are limiting normals. Set prox-regularity at
$x_0$ requires a uniform local projection property near each limiting normal
at $x_0$ \cite[Definition~1.1 and Proposition~1.2]{PRT}. In particular, every
limiting normal at $x_0$ must itself be proximal there. For closed sets in finite
dimension, set prox-regularity at $x_0$ is equivalent to single-valued nearest
projection on some neighborhood of $x_0$ \cite[Theorem~1.3(a),(k)]{PRT}.
These facts apply in any fixed inner product, by a linear isometry to Euclidean
space. Locally closed sets may be clipped inside a closed ball, with the
base point in its interior.

For $t>0$ take the stable path
\begin{equation}\label{eq:path}
 p_t(z)=z^2+t^2z+t^2+it^3=(z+it)(z+t^2-it),
 \qquad c_t=(t^2,0,t^2,t^3).
\end{equation}
The root $\lambda=-it$ is simple, with the other root strictly to its left.
Thus $\alpha$ is smooth near $c_t$. Implicit differentiation in all complex
coefficient directions gives
\[
 d\lambda=-\frac{\lambda\,da+db}{t^2-2it},\qquad
 \nabla\alpha(c_t)=\gamma_t n_t,\qquad
 \gamma_t=\frac{2}{t(t^2+4)},\quad
 n_t=(-t,-t^2/2,-t/2,1).
\]
Here $n_t$ is a real support covector. In metric $G$ its representing normal
vector is $G^{-1}n_t$.

At each fixed $t$, Taylor's theorem for the local smooth constraint
$\alpha\leq0$ yields
$n_t^T(c-c_t)\leq K_t\norm{c-c_t}_G^2$ for nearby stable $c$.
Expanding squared distances shows that, for a sufficiently small positive
$\tau$, $c_t$ is the unique closest such point to $c_t+\tau G^{-1}n_t$.
Choose the step smaller also than half the radius of that local metric ball;
the triangle inequality excludes all points outside the ball. Consequently
$G^{-1}n_t$ is an actual proximal normal to $S$, without requiring a uniform
$K_t$ or a uniform root gap.

Let $n_0=(0,0,0,1)$. Since $c_t\to c_0$ and $n_t\to n_0$,
$G^{-1}n_0$ is a limiting normal at $c_0$. However,
\begin{equation}\label{eq:obstruction}
 n_0^Tc_t=t^3,\qquad
 \norm{c_t}^2=2t^4+t^6.
\end{equation}
More generally, writing $e=(1,0,1,0)$ and $f=(0,0,0,1)$,
\[
 \norm{c_t}_G^2=t^4(e^TGe+2t e^TGf+t^2f^TGf),\qquad e^TGe>0.
\]
For every $\tau>0$,
\[
 \norm{c_t-\tau G^{-1}n_0}_G^2
       -\norm{\tau G^{-1}n_0}_G^2
       =\norm{c_t}_G^2-2\tau t^3<0
\]
when $t$ is sufficiently small. The competitors converge to the base point.
Thus $G^{-1}n_0$ is not even a local proximal normal at $c_0$, disproving
prox-regularity for every fixed $G$.

The continuous quadratic formula
\begin{equation}\label{eq:alpha}
 \alpha(a,b)=-\frac{\Rea a}{2}
  +\frac12\sqrt{\frac{|a^2-4b|+\Rea(a^2-4b)}{2}}
\end{equation}
shows that $S$ is closed. It is nonempty and all distances to it are attained:
a minimizing sequence is bounded by comparison with $c_0$, and has a
convergent subsequence. The projection characterization therefore gives
non-single-valued projections in every neighborhood of $c_0$, with at least
two distinct minimizers. Their inputs are unstable, since a stable input has
only itself as a distance-zero nearest point. This proves the monic assertion.

\begin{remark}[Boundary chart and real coefficients]\label{rem:chart}
Writing $a=x+iy,b=u+iv$ and substituting $z=w-iy/2$ gives
$w^2+xw+U+iV$, where $U=u+y^2/4$ and $V=v-xy/2$.
The quadratic formula gives the exact closed stability chart
\[
 x\geq0,\qquad U\geq0,\qquad V^2\leq x^2U.
\]
Indeed its largest root real part is
$(-x+\sqrt{(|x^2-4U-4iV|+x^2-4U)/2})/2$.
Stability is equivalent to $x\geq0$ and
$|x^2-4U-4iV|\leq x^2+4U$. Squaring with its right side nonnegative gives
precisely the displayed inequalities, including $x=0$ and $U=0$.
The path \eqref{eq:path} lies on the upper smooth
sheet $V=x\sqrt U$. This is the positive umbrella solid discussed above.
For real monic quadratics the stability set instead is the convex orthant
$a\geq0,b\geq0$; projection is unique for every fixed inner product. The
imaginary constant coefficient in \eqref{eq:path} is essential to this example.
\end{remark}

\section{Full degree and the epigraph}
In coordinate order
$(\Rea\ell,\Ima\ell,\Rea a,\Ima a,\Rea b,\Ima b)$,
direct differentiation at the same monic points gives
\[
 d\lambda=-\frac{\lambda^2\,d\ell+\lambda\,da+db}{t^2-2it},\qquad
 \nabla\alpha(1,t^2,t^2+it^3)=\gamma_t N_t,
\]
\[
 N_t=(t^3/2,-t^2,-t,-t^2/2,-t/2,1)
       \longrightarrow N_0=(0,0,0,0,0,1).
\]
Both leading-coefficient directions are included. The earlier argument makes
$G^{-1}N_t$ proximal locally for each full coefficient metric $G$.
The displacement from $c_{\rm full}$ has
zero leading coordinates, support $t^3$, and a squared metric norm of order
$t^4$ with a positive leading coefficient. Hence the limiting normal
$G^{-1}N_0$ is not proximal.

The exact-degree stable set is not globally closed, so attainment requires
localization. On $\ell\ne0$ the abscissa equals $\alpha(a/\ell,b/\ell)$ and
is continuous. Choose a closed metric ball $B$ of radius $R$ about
$c_{\rm full}$ with $\ell\ne0$ throughout. Then
$C=S_{\rm full}\cap B$ is nonempty and compact. Its local normals along
the path are unchanged, so $C$ fails prox-regularity at the base. For an
input $q$ with $\norm{q-c_{\rm full}}_G<R/3$, every stable point outside
$B$ is more than $2R/3$ from $q$, while the stable base is less than $R/3$
away. Boundary points of $B$ cannot minimize either. Thus the global nearest
points in $C$ and $S_{\rm full}$ agree. Applying the closed-set projection
characterization to $C$ proves the full-degree assertion of
Theorem~\ref{thm:main}, including global attainment and unstable inputs.

At $(c_t,0)$ the epigraph has smooth constraint $\alpha(c)-r\leq0$.
Its outward proximal normal covectors can be scaled to
\[
 v_t=(n_t,-1/\gamma_t)\longrightarrow(n_0,0).
\]
For full degree replace $n_t$ by $N_t$.
For any fixed positive-definite coefficient/height metric $H$, their
representing proximal normal vectors are $H^{-1}v_t$. Comparison with the
epigraph points $(c_t,0)$ again gives support $t^3$ and squared distance of
order $t^4$. The limit is not proximal at height zero, proving the epigraph
assertion. The monic epigraph is closed; in full degree a small closed ball
inside $\ell\ne0$ supplies local closedness. The limiting support covector is
horizontal. Its representing vector is horizontal for a coefficient/height
product metric, and can tilt when the metric mixes those coordinates.

\section{Why the function condition is different}
Function prox-regularity at a point relative to a finite limiting subgradient
$v_0$ localizes the subgradients near $v_0$ as well as the base points and
function values. Its epigraph correspondence is at the normal $(v_0,-1)$
\cite[Definition~1.1 and p.~205]{PR}. The preceding horizontal normal does
not test that condition: the unscaled gradients along \eqref{eq:path}
diverge. The distinction is substantive, as the next proposition shows.

\begin{proposition}\label{prop:function}
The monic complex quadratic abscissa \eqref{eq:alpha} is function prox-regular
at $(0,0)$ relative to every finite limiting subgradient, in the real Euclidean
coefficient metric. The constants may depend on the specified subgradient.
\end{proposition}
\begin{proof}
Put $D=a^2-4b$, $L=(-\infty,0]\subset\C$, and
$\Phi(D)=\sqrt{(|D|+\Rea D)/2}$. Real pairings on $\C$ are understood.
For $D\notin L$, $\Phi$ is smooth and
$\norm{\nabla\Phi(D)}=1/(2\sqrt{|D|})$.
At $D=-s<0$, rationalizing the square root gives
\[
 \Phi(D+h)=\frac{|\Ima h|}{2\sqrt{s}}+O(|h|^2),
\]
so its regular, proximal, and limiting subdifferentials are the imaginary
interval $\{iq:|q|\leq1/(2\sqrt{s})\}$. The limiting statement follows
also by the two off-ray gradient limits at that fixed point.

For $|D|\leq\delta$ one has
\begin{equation}\label{eq:distance}
 \Phi(D)\geq\frac{\dist(D,L)}{2\sqrt\delta}.
\end{equation}
For $\Rea D\geq0$, use $\dist(D,L)=|D|$ and
$\Phi(D)\geq\sqrt{|D|/2}$. For $\Rea D<0$, use
\[
 \Phi(D)=\frac{|\Ima D|}{\sqrt{2(|D|-\Rea D)}}.
\]
At zero every finite $q$ in the convex normal cone
$N_L(0)=\{q:\Rea q\geq0\}$ has local linear lower support by
\eqref{eq:distance}, after choosing $\delta$ sufficiently small.
Conversely, increments on the negative real ray force $\Rea q\geq0$.
Off-ray gradients diverge in norm as $D\to0$, while bounded on-ray
subgradients belong to $N_L(0)$. Hence all three subdifferentials at zero
equal $N_L(0)$, with the limiting conclusion valid also when defined using
limits of proximal subgradients.

The global real smooth diffeomorphism $(a,b)\mapsto(a,a^2-4b)$, with inverse
$b=(a^2-D)/4$, transports these subdifferentials exactly. First-order lower
supports transform by the Jacobian; quadratic proximal supports transform
by Taylor's theorem, and limits follow by continuity of the inverse
Jacobian. In particular any limiting subgradient $v$ of $\alpha$ has
\begin{equation}\label{eq:subgrad}
 v=-\tfrac12e_{\Rea a}+\tfrac12 J_D(a,b)^Tq,
 \qquad v_b=-2q,\qquad q\in\partial\Phi(D).
\end{equation}
This is a direct special case of the standard smooth composition calculus
\cite[Theorem~3.1]{PR}.

Fix a finite $v_0\in\partial\alpha(0,0)$. In a fixed neighborhood of $v_0$,
\eqref{eq:subgrad} bounds $|q|\leq M$, for some $M\geq1$.
Restrict the coefficient ball so $|D|<\delta=1/(16M^2)$ throughout.
Every nearby base $x$ with a subgradient near $v_0$ must have $D_x\in L$:
an off-ray gradient would have norm greater than $2M$.
Then $\Phi(D_x)=0$ and $q\in N_L(D_x)$.
For any comparison point $y$ in that ball and its nearest-ray point $z$,
convex normal support and \eqref{eq:distance} give
\[
 q\mathbin{\cdot}(D_y-D_x)
 \leq M\dist(D_y,L)\leq\Phi(D_y).
\]
Writing $\Delta=y-x$ and using the exact identity
$D_y-D_x=J_D(x)\Delta+(\Delta a)^2$ now yields
\[
 \alpha(y)\geq\alpha(x)+v\mathbin{\cdot}\Delta
       +\tfrac12q\mathbin{\cdot}(\Delta a)^2
 \geq\alpha(x)+v\mathbin{\cdot}\Delta-\tfrac M2\norm\Delta^2.
\]
With curvature constant $M+1$ the inequality is strict for $y\ne x$, as in
\cite[Definition~1.1]{PR}. It holds without needing a restriction on the
base function value, hence in particular on the required function-attentive
neighborhood. Formula \eqref{eq:alpha} supplies continuity. This proves the
proposition.
\end{proof}

Cross \cite[Chapter~5]{Cross} studies the finite-subgradient function
condition and gives a negative example for a nonlinear two-real-parameter
matrix family (Example~5.38). Its coefficient map is
$(x_1,x_2)\mapsto(a=0,b=-x_1^2+x_2^2)$, whose derivative at the origin
vanishes. That example concerns the parameter metric and does not identify
distinct nearest complex coefficient polynomials. His real-quadratic
pointwise proximal-subgradient calculation on pp.~131--133 is also a
different statement. Neither such parameter results nor partial smoothness
replace the ambient support and attainment arguments above.

\section*{Reproducibility and disclosure}
The verification package contains this source, a standard-library exact
rational verifier for the root, normal, and support identities, deterministic
high-precision boundary controls for Proposition~\ref{prop:function}, and
an authored source/priority memorandum. Finite checks support reproducibility;
the limiting, universal-metric, and global projection conclusions follow
from the proofs. No explicit unstable input together with its two minimizers
is claimed, and no all-degree or real-only classification is supplied.

AI tools were used extensively in deriving, repairing, checking, literature
auditing, and writing this work, including independent adversarial AI-agent
reviews. These reviews are not human peer review. This is an unrefereed
preprint, without formal proof-assistant certification; the named author
is responsible for the claims and citations. No outside individual was
contacted during the research. The paper and accompanying original materials
are licensed under Creative Commons Attribution 4.0 International.

\begin{thebibliography}{10}
\bibitem{Bottema} O. Bottema, \emph{The Routh-Hurwitz condition for the
biquadratic equation}, Indagationes Mathematicae 18 (1956), 403--406.
\href{https://doi.org/10.1016/S1385-7258(56)50054-6}{doi:10.1016/S1385-7258(56)50054-6}.
\bibitem{Lewis2003} A.~S. Lewis, \emph{The mathematics of eigenvalue
optimization}, Mathematical Programming 97 (2003), 155--176, especially
\S8, pp.~166--168.
\href{https://people.orie.cornell.edu/aslewis/publications/03-mathematics.pdf}{Author-hosted paper}.
\bibitem{OWR} A.~S. Lewis, \emph{Local Structure and Algorithms in Nonsmooth
Optimization} (joint work with J.~V. Burke and M.~L. Overton), in
\emph{Optimization and Applications}, Oberwolfach Reports 2 (2005),
104--106, Question~1 on p.~106.
\href{https://doi.org/10.4171/OWR/2005/02}{doi:10.4171/OWR/2005/02}.
\bibitem{BO} J.~V. Burke and M.~L. Overton, \emph{Variational analysis of
the abscissa mapping for polynomials}, SIAM Journal on Control and
Optimization 39 (2001), 1651--1676.
\href{https://sites.math.washington.edu/~burke/papers/reprints/37-va-abs-poly.pdf}{Author-hosted paper}.
\bibitem{BLO} J.~V. Burke, A.~S. Lewis, and M.~L. Overton,
\emph{Variational analysis of functions of the roots of polynomials},
Mathematical Programming 104 (2005), 263--292.
\href{https://doi.org/10.1007/s10107-005-0616-1}{doi:10.1007/s10107-005-0616-1}.
\bibitem{PRT} R.~A. Poliquin, R.~T. Rockafellar, and L. Thibault,
\emph{Local differentiability of distance functions}, Transactions of the
American Mathematical Society 352 (2000), 5231--5249.
\href{https://sites.math.washington.edu/~rtr/papers/rtr170-LocalDiffDistance.pdf}{Author-hosted paper}.
\bibitem{KO} O.~N. Kirillov and M.~L. Overton, \emph{Robust stability at
the Swallowtail singularity}, Frontiers in Physics 1 (2013), Article~24.
\href{https://doi.org/10.3389/fphy.2013.00024}{doi:10.3389/fphy.2013.00024}.
\bibitem{PR} R.~A. Poliquin and R.~T. Rockafellar, \emph{A calculus of
prox-regularity}, Journal of Convex Analysis 17 (2010), 203--210.
\href{https://www.heldermann-verlag.de/jca/jca17/jca0849_b.pdf}{Publisher paper}.
\bibitem{Cross} J.~A. Cross, \emph{Spectral Abscissa Optimization using
Polynomial Stability Conditions}, Ph.D. dissertation, University of
Washington, 2010, Chapter~5, pp.~80--133.
\href{https://www.proquest.com/docview/288328008}{Open-access primary record}.
\bibitem{ML} R.~L. Moses and D. Liu, \emph{Determining the closest stable
polynomial to an unstable one}, IEEE Transactions on Signal Processing 39
(1991), 901--906.
\href{https://doi.org/10.1109/78.80912}{doi:10.1109/78.80912}.
\bibitem{Orban} F.-X. Orban de Xivry, \emph{Nearest Stable System},
doctoral thesis, Universit\'e catholique de Louvain, 2013,
Chapters~2 and~4.
\href{https://research.dial.uclouvain.be/entities/publication/650a4c86-3c4c-439e-9a8f-54c4d42eab7f}{Open-access primary record}.
\end{thebibliography}
\end{document}
